\hypertarget{r29-am2-am2-forward-hnm-four-site-quantitative-stability-certificate}{%
\section{HNM AM2: four-site quantitative stability certificate\workstar}\label{r29:am2:r29-am2-am2-forward-hnm-four-site-quantitative-stability-certificate}}
% BEGIN CURRENT REGISTRY NOTE
\paragraph{Statement alias HNM-T-AM2.}\label{stmt:AM2} This identifier denotes the scoped mathematical claim admitted by this chapter\textquotesingle s gate. It adds no assumption and establishes no literature-priority claim.
\begin{quote}\small\textbf{Current extension: HNM-C-AQ1.} The historical limitation below remains the scope of this original result. The later, separately gated statement in Section~\ref{stmt:AQ1} records the extension: At both signs of |tau|<=1/100000000, actual centered full-Z3 whole-star boxes have a subsequence defining a compatible locally normal gauge-invariant stationary state on the norm closure of full bounded local algebras. The actual reference reset bound is 56|tau||R|; compact-resolvent cutoffs prove local trace-norm precompactness. The matched unbounded-onsite dynamics theorem and separate local-normality argument supply strongly continuous GNS evolution with a nonnegative self-adjoint physical energy generator. The forward physical cyclic restriction is reducing; full joint-fixed completion and its positive gap remain for AQ2. Its own model and state assumptions remain required.\end{quote}
\begin{quote}\small\textbf{Current extension: HNM-C-AQ2.} The historical limitation below remains the scope of this original result. The later, separately gated statement in Section~\ref{stmt:AQ2} records the extension: For the actual AQ1 centered full-Z3 subsequential state at both signs of |tau|<=1/100000000, the complete original-endpoint gauge-fixed space equals the invariant-local cyclic completion and reduces the physical energy generator. Centered Fourier tests including zero give H\_phys >= (alpha/16)(I-P\_Omega) and a simple vacuum in this representation. The full-GNS strengthening explicitly uses AM2's separately reviewed full-Hilbert finite gap. For the original xz Wilson in the same state, the actual 48-link/36-endpoint cover has seven incident stars, local reference energy <=98|tau| and variance >=61999/250000>1/5. The reverse proof separately sharpens that floor to 3099951/12500000. Its own model and state assumptions remain required.\end{quote}
% END CURRENT REGISTRY NOTE

\begin{quote}\small \textbf{Admission: accepted within scope.} For every nonempty finite I1 complete-factor volume with the actual selected-strip reference and full whole-star omitted interaction, both signs of \(|\tau|\le 1/100000000\) admit a unique ground and a physical gap at least \(\alpha/16\). The proof derives the four-site creation majorant \(G(t)=16\) exp(8t)(1+10t), full scalar fixed point, centered excited-sector exclusion and removal of onsite representation cutoffs. The forward full-Hilbert-space strengthening is separately proved and reviewed.\par The typeset body below preserves the source-bound forward derivation. Its prospective language is historical; this boxed gate records the final reviewed verdict. The separately authored reverse and skeptical records remain linked. The common admission and attributed refinements are determined by the gate, not by a renamed equation.\end{quote}
Human project author: \textbf{Hruday N M (BUNZEEY)}. The advisor/panel supplied the proposed route and constants. This independent forward derivation checks the entire implication before reading current reverse or skeptic results. Creation expansions are established methods; scientific priority of this application is unverified.

\textbf{Proposed admission:} for every finite I1 complete-factor volume, with its actual selected reference, whole-star omitted interactions and both signs of \(|\tau|\le 10^{-8}\), the full untruncated Hamiltonian has a unique ground and centered gap at least \(\alpha/16\). Its physical restriction has the same ground and this lower gap. Every nonempty complete-factor volume has nonzero physical excited vectors; the empty volume has none and its gap exclusion is vacuous. This is an alternative explicit finite-volume theorem, not an evaluation of Yarotsky's named constants or a continuum result.

\hypertarget{r29-am2-actual-model-and-available-inverse}{%
\subsection{Actual model and available inverse}\label{r29:am2:r29-am2-actual-model-and-available-inverse}}

Write the normalized operator on a finite set \(\Lambda\) of complete 24-link factors as

\begin{equation}
H=H_0+V,\quad H_0=\sum_{x\in\Lambda}h_x,\quad
h_x\Omega_x=0,\quad h_x\ge Q_x=I-|\Omega_x\rangle\langle\Omega_x|,
\quad V=\sum_{b:b+S\subset\Lambda}\phi_b .
\tag{HNM-AM2.1}
\label{r29:am2:eq1}
\end{equation}

All 21 omitted faces per retained anchor are present, \(S=\{0,e_x,e_y,e_z\}\), and \(\|\phi_b\|=7|\tau|\). The indexed per-site sum is

\begin{equation}
J:=\max_u\sum_{X\ni u}\|V_X\|\le28|\tau|\le J_0={7\over25000000}.
\tag{HNM-AM2.2}
\label{r29:am2:eq2}
\end{equation}

Here each retained X is its complete four-site star, including incoming anchors. The interaction is bounded in each finite volume, though not uniformly bounded as a global extensive operator. All onsite domains and scalar selected ground subtractions are inherited from I1. No interaction term is deleted.

For a nonempty excitation set M, put \(P_M=(\bigotimes_{x\in M}Q_x)(\bigotimes_{y\notin M}P_y)\) and identify its range with the excited factor tensor product. On it \(H_M=\sum_{x\in M}h_x\ge |M|\). This is a full-Hilbert-space inequality and needs no extra physical four-link floor. For real \(|z|<1/2\),

\begin{equation}
\|H_M^{-1}\|\le|M|^{-1},\qquad
\|(H_M-z)^{-1}\|\le2|M|^{-1}.
\tag{HNM-AM2.3}
\label{r29:am2:eq3}
\end{equation}

First impose finite-rank full spectral cutoffs of each \(h_x\) containing its vacuum. The proof below is finite dimensional and constants do not depend on these ranks. Cutoff removal is proved in Section6.

\hypertarget{r29-am2-creation-algebra-and-complete-anchored-estimate}{%
\subsection{Creation algebra and complete anchored estimate}\label{r29:am2:r29-am2-creation-algebra-and-complete-anchored-estimate}}

For a vector \(c_I\in\bigotimes_{x\in I}Q_x\mathcal H_x\), define the identity-extended creation \(\widehat{c_I}=|c_I\rangle\langle\Omega_I|\otimes I_{I^c}\). Two creations commute: disjoint supports commute by tensor factorization, and overlapping supports have product zero in both orders. The latter follows because the vacuum bra annihilates an excited ket at any common factor, first for simple tensors and then by finite-dimensional linearity. Set

\(C=\sum_{I\ne\varnothing}\widehat{c_I}\), \(\|c\|_a=\max_u\sum_{I\ni u}\|c_I\|\).

This is a norm on the finite-dimensional direct sum of coefficient spaces. For collections \(c_1,...,c_k\), define

\begin{equation}
(L_k(c_1,\ldots,c_k))_M
=H_M^{-1}P_M\operatorname{ad}_{C_1}\cdots\operatorname{ad}_{C_k}(V)\Omega_0,
\quad\Omega_0=\bigotimes_x\Omega_x .
\tag{HNM-AM2.4}
\label{r29:am2:eq4}
\end{equation}

We derive, rather than copy, the proposed estimate

\begin{equation}
\|L_k(c_1,\ldots,c_k)\|_a
\le J L_k^{\rm num}\prod_j\|c_j\|_a,
\qquad L_k^{\rm num}=16\,8^k(1+5k/4).
\tag{HNM-AM2.5}
\label{r29:am2:eq5}
\end{equation}

Fix an indexed interaction \(V_X\), \(|X|\le p=4\), and creation supports \(I_j\). Every \(I_j\) must meet X: commuting adjoint actions allow a disjoint creation to be moved innermost, where it commutes with \(V_X\). Write \(N=\bigcup_j I_j\). In any nonzero term, outside X a singly touched factor is excited and a multiply touched factor annihilates the term. Therefore \(N\setminus X\subset M\subset N\cup X\). There are at most \(2^p=16\) possible output support sets M, independent of Hilbert dimensions, and \(|I_j|\le |M|+p\). The commutator expansion has \(2^k\) products; each projected product norm is at most \(\|V_X\|\prod_j\|c_{j,I_j}\|\).

For a fixed root u in M, first count terms with u in X. Use inverse bound one, sum X containing u at cost J, and for every collection use \(\sum_{I:I\cap X\ne\varnothing}\|c_I\|\le p\|c\|_a\). This contributes \(2^p(2p)^k J\prod_j\|c_j\|_a\).

Otherwise u lies in at least one \(I_l\); overcount all k choices of l. Since M is nonempty,

\(1/|M| \le  (p+1)/|I_l|\).

Sum interactions meeting \(I_l\) at cost \(J|I_l|\), canceling that denominator; sum the remaining collections meeting X at cost \(p^(k-1)\), and sum \(I_l\) containing u with its anchored norm. This contributes \(2^p(2p)^k k(p+1)/p\) times the same J product. Adding the two root placements and setting \(p=4\) gives (HNM-AM2.5). All sums are finite; repeated interaction supports and repeated creation histories are counted, not canceled to improve the bound.

The nested commutator is zero for \(k>2p=8\). In each expanded product, one side of \(V_X\) would contain at least five creations all meeting the four-site X. Two share a site, so that product is zero. An exact four-qubit algebra fixture gives \(\operatorname{ad}_C^8(V)\Omega=8!|1111\rangle\) and \(\operatorname{ad}_C^9(V)=0\); thus importing a support-three termination at order six would be wrong. This fixture audits the algebra, not a rotor truncation or a numerical spectrum.

\hypertarget{r29-am2-fixed-point-scalar-energy-and-all-orders}{%
\subsection{Fixed point, scalar energy and all orders}\label{r29:am2:r29-am2-fixed-point-scalar-energy-and-all-orders}}

Because creators commute, \([C,H_0]=-\sum_I\widehat{H_Ic_I}\) and its next creation commutator vanishes. Thus

\begin{equation}
e^C(H_0+sV)e^{-C}=H_0-\sum_I\widehat{H_Ic_I}+s e^CVe^{-C}.
\tag{HNM-AM2.6}
\label{r29:am2:eq6}
\end{equation}

All exponential products are finite polynomials: C is nilpotent since nonzero creator products have disjoint nonempty supports. For real \(|s|\le 1\), \(\psi(s)=e^{-C(s)}\Omega_0\) is an eigenvector precisely when

\begin{equation}
c=s\sum_{k=0}^{8}{L_k(c,\ldots,c)\over k!},\qquad
E(s)=s\langle\Omega_0,e^{C(s)}Ve^{-C(s)}\Omega_0\rangle .
\tag{HNM-AM2.7}
\label{r29:am2:eq7}
\end{equation}

The vacuum coefficient of \(\psi\) is one, so it never vanishes. The scalar equation is retained even when the original reference vacuum mean of V is zero. All diagonal pieces of V remain inside every commutator and the scalar; this is not an offdiagonal-only model.

The positive all-order majorant and derivative are

\begin{equation}
G(t)=\sum_{k\ge0}{L_k^{\rm num}t^k\over k!}
=16e^{8t}(1+10t),\quad
G'(t)=16e^{8t}(18+80t).
\tag{HNM-AM2.8}
\label{r29:am2:eq8}
\end{equation}

For \(R=1/64\), \(\exp(1/8)<8/7\), by comparison of its positive power series with the geometric series (strict already at second order). Hence

\begin{equation}
\begin{aligned}G(R)&<148/7,\qquad G'(R)<352,\\ J_0G(R)&<148/25000000<R,\\ J_0G'(R)&<77/781250,\qquad 2J_0G'(R)<77/390625<1.\end{aligned}\tag{HNM-AM2.9}
\label{r29:am2:eq9}
\end{equation}

The multilinear bound proves the ball self-map. Telescoping each k-linear term between two coefficient collections in the radius-R ball yields the Lipschitz bound \(J G'(R)\), with all k possible replacements counted. Banach's theorem supplies a unique coefficient solution there. Subtracting the fixed-point equations and absorbing the common Lipschitz factor gives continuity in s. Therefore E(s) is a continuous real eigenvalue branch. At \(J=0\), \(c=0\) and \(E=0\) exactly.

\hypertarget{r29-am2-excited-sector-exclusion-construction-alone-is-not-enough}{%
\subsection{Excited-sector exclusion: construction alone is not enough}\label{r29:am2:r29-am2-excited-sector-exclusion-construction-alone-is-not-enough}}

Suppose a second eigenvector phi of \(H_0+sV\) has eigenvalue \texttt{E(s)+z}, real \(|z|<1/2\), and is independent of psi. Decompose

\(e^{C(s)}\phi=b_0\Omega_0+\sum_{M\ne\varnothing}b_M\), and set \(B=\sum_M\widehat{b_M}\).

Then \(\phi=(b_0+B)\psi\) because creations commute, and b is nonzero. Subtract the eigenvector equations, conjugate by \(e^C\), and project each nonempty sector. Creation commutation cancels the creation term in (HNM-AM2.6), giving

\begin{equation}
b_M=s(H_M-z)^{-1}P_M[B,e^{C(s)}Ve^{-C(s)}]\Omega_0.
\tag{HNM-AM2.10}
\label{r29:am2:eq10}
\end{equation}

The scalar equation may additionally constrain \(b_0\), but is unnecessary for excluding \(b!=0\). Expanding the remaining commutator keeps all terms through total creation order eight. The shifted inverse multiplies (HNM-AM2.5) by at most two. Treating b as one input and summing the other orders gives

\(\|b\|_a\le 2JG'(R)\|b\|_a<\|b\|_a\), a contradiction.

Thus the branch is simple and no other eigenvalue lies within one half. At \(s=0\) it is the unique lowest eigenvalue; continuous finite-dimensional eigenvalues cannot cross this isolated branch. It remains the ground for every real \(|s|\le 1\), with full-space gap at least one half. The argument includes \(\tau\) of either sign and equality at the contracted \(\tau\) cap. It does not interpret rejection beyond the chosen cap as actual gap failure.

\hypertarget{r29-am2-gauge-and-physical-sector}{%
\subsection{Gauge and physical sector}\label{r29:am2:r29-am2-gauge-and-physical-sector}}

Every onsite selected-strip/free-link operator commutes with the original gauge actions at all owned-link endpoints, including heads outside the tail volume. Its unique positive ground is fixed by these actions. Hence \(P_x,Q_x\), every spectral cutoff and \(H_M\) commute with gauge actions. The fixed-point map preserves invariant collections; uniqueness gives invariant C and psi. Equivalently the unique full ground is fixed under the finite product of endpoint SU(2) groups. The physical subspace reduces the full Hamiltonian, so the full-space gap restricts to it with the same vacuum.

For a nonempty complete-factor volume, choose a selected square in any factor. The full finite-volume ground is smooth and strictly positive on the connected compact product of link groups, including both signs of its bounded real potential. The gauge-invariant vector \((W_p-\langle W_p\rangle)\Omega_H\) is nonzero, is in the operator domain and is orthogonal to the ground. This proves a genuine nonzero physical excited sector. The empty volume is one dimensional and is treated separately.

\hypertarget{r29-am2-untruncated-compact-domain-passage}{%
\subsection{Untruncated compact-domain passage}\label{r29:am2:r29-am2-untruncated-compact-domain-passage}}

Each actual \(h_x\) is an elliptic Casimir sum plus bounded smooth selected potential and a scalar shift on the compact manifold \(SU(2)^{24}\). It is self-adjoint on the Casimir domain and has compact resolvent. Choose increasing full spectral projections \(Q_{x,N}\) containing its ground and tending strongly to the identity. They commute with gauge actions. The product cutoff \(Q_N\) commutes with \(H_0\), and compressed interactions retain support and do not increase J. The preceding estimates are uniform in N.

Finite tensor products of onsite eigenvectors form a form core for \(H_0\): spectral tails of \(\|u\|^2+\langle u,H_0u\rangle\) tend to zero. Since V is bounded in each fixed finite volume, the same union is a form core for \(H_0+sV\), after a harmless scalar lower-bound shift. Compact resolvent persists. By min--max, the first two cutoff eigenvalues converge downward to the first two full eigenvalues: form approximation to their finite-dimensional eigenspaces proves the upper limit, while the variational principle gives the lower limit. The uniform cutoff inequality therefore gives \(E_1-E_0\ge 1/2\) for the untruncated operator. The same core argument applies on the physical subspace because \(Q_N\) commutes with gauge actions; its nonzero sector was established separately.

No infinite-volume bounded global creator, similarity transform or unitary is asserted. The normalized physical energy is restored using the original scalar relation: after actual ground subtraction,

\begin{equation}
\Delta_{\rm physical}\ge{\alpha\over8}\,{1\over2}={\alpha\over16},
\qquad\Delta/E_\star\ge{\alpha\over16E_\star},
\quad\hbox{frequency threshold }\alpha/(16\hbar).
\tag{HNM-AM2.11}
\label{r29:am2:eq11}
\end{equation}

\hypertarget{r29-am2-evidence-and-scope}{%
\subsection{Evidence and scope}\label{r29:am2:r29-am2-evidence-and-scope}}

The exact checker verifies the anchored combinatorial constants, complete four-site termination example, certified exponential enclosure, radius inequalities, shifted-inverse bound, zero/equality cases, scalar/diagonal controls and the I1 J conversion. These audit the analytic proof; passing a polynomial or a finite cutoff alone is not this theorem. The admitted result, if independent review confirms the argument, is a conservative uniform \textbf{finite-volume} radius for the complete original I1 family. It does not evaluate Yarotsky's c1,c2, select an infinite-volume state, prove AN's boundary comparison or repair AL's weak-coupling obstruction. Subsequent infinite-volume or continuum claims need their own complete arguments.

Reproduce from the repository root:
\begin{verbatim}
python -B research/round29/forward/am2/check.py \
  --output /absolute/new/output
\end{verbatim}

\subsection*{Final admitted limits and evidence}
\begin{itemize}
\item This is a deliberately conservative finite-volume theorem uniform over the declared volumes, not an infinite-volume state-identification theorem.
\item The proof does not evaluate Yarotsky's named c1/c2 or HTW's boundary constants; inherited infinite-state claims retain their own restrictions.
\item The empty volume has only the vacuum and spectral exclusion is vacuous; nonzero physical excitations are asserted for nonempty complete-factor volumes.
\item Fixed positive physical scales, actual onsite reference and complete interaction grouping are essential. No weak-bare-coupling continuum bridge or continuum mass-gap construction follows.
\item The candidate constants were shared panel proposals and the creation method is established mathematics; scientific priority remains unverified.
\end{itemize}
\repo{research/round29/advisor/am2-gate.json}\par\repo{research/round29/reverse/am2/report.md}\par\repo{research/round29/skeptic/am2.md}

\emph{Sequencing note.} Prospective statements in this frozen derivation describe the moment of its admission. The final Round29 roadmap below governs subsequent work.
